\documentclass[12pt]{article}
\usepackage[utf8]{inputenc}
\usepackage[T1]{fontenc}
\usepackage{hyperref}
\title{Wanggamala Language Resources for K-12 Curriculum Development}
\author{Compiled by Claude Sonnet 5.5 (claude-sonnet-5-5) for Andr\'as Kornai}
\date{}
\begin{document}
\maketitle

\section{Introduction}
Wanggamala (ISO 639-3: wnm; Glottocode wang1289) is an extinct Pama-Nyungan language of Australia. Sources place it in the Northern Territory around the Hay River, east of Alice Springs. Parts of central west Queensland are also associated with it. Wikipedia reports one remaining speaker as of 2003. Variant spellings include Wangkamanha, Wongkamala and Wonkamudla. The Endangered Languages Project still lists it, which allows for revitalisation by people who identify with the language.

\section{Primary Educational Resources}
No primers, textbooks or school programmes for Wanggamala were found in the searches.

\section{Linguistic and Reference Materials}
\begin{itemize}
\item Wikipedia entry on the Wanggamala language: \url{https://en.wikipedia.org/wiki/Wanggamala_language}
\item Glottolog record for Wanggamala: \url{https://glottolog.org/resource/languoid/id/wang1289}
\item PHOIBLE inventory page for the language code: \url{https://phoible.org/languages/wang1289}
\end{itemize}

\section{Religious and Community Resources}
None were found in the searches.

\section{Digital and Audio Resources}
The Endangered Languages Project catalogue entry is at \url{https://endangeredlanguages.com/lang/3330}. No audio, video or corpus resources were found in the searches.

\section{Curriculum Development Context}
Little documentation appears in the sources found. A heritage curriculum would have to start from archival and linguistic materials, which these searches did not locate. It should be developed together with descendant communities. The most reasonable first step is to approach the catalogue and archives for existing wordlists and recordings.
\end{document}